\documentclass[12pt]{article}
\usepackage[a3paper, margin=1in]{geometry}
\usepackage{amsmath, amssymb, amsthm, ulem, booktabs}

\begin{document}

\section*{Domácí úkol 9.2}

Vypracoval: Daniel \textit{``Randál"} Ransdorf \hfill Podpis: \rule{4cm}{0.4pt}

\begin{enumerate}
  
  \item Množina X je množinou všech řešení rovnice $2x_1-3x_2+4x_3+x_4=0$, rovnici zapišme do matice a vyřešme:
    \[
      \left(\begin{array}{cccc|c}
        2 & -3 & 4 & 1 & 0
      \end{array}\right)
      \overset{parametry}{\sim}
      \left(\begin{array}{cccc|c}
        2 & -3 & 4 & 1 & 0 \\ \hline
        0 & 1 & 0 & 0 & 2t_1 \\
        0 & 0 & 1 & 0 & t_2 \\
        0 & 0 & 0 & 1 & 2t_3
      \end{array}\right)
      \sim
      \left(\begin{array}{cccc|ccc}
        2 & 0 & 0 & 0 & 6t_1&-4t_2&-2t_3 \\ \hline
        0 & 1 & 0 & 0 & 2t_1 \\
        0 & 0 & 1 & 0 & &t_2 \\
        0 & 0 & 0 & 1 & &&2t_3
      \end{array}\right)
      \sim
      \left(\begin{array}{cccc|ccc}
        1 & 0 & 0 & 0 & 3t_1&-2t_2&-t_3 \\ \hline
        0 & 1 & 0 & 0 & 2t_1 \\
        0 & 0 & 1 & 0 & &t_2 \\
        0 & 0 & 0 & 1 & &&2t_3
      \end{array}\right)
    \]

    Z tohoto lze vidět:
    \[
    X = Span\left\{
        \left(\begin{array}{c} 3\\2\\0\\0 \end{array}\right),
        \left(\begin{array}{c} -2\\0\\1\\0 \end{array}\right),
        \left(\begin{array}{c} -1\\0\\0\\2 \end{array}\right)
      \right\}
    \]

    Zadefinujme si matice $A_X,A_Y$, které mají vektory ze $Span$ jako sloupce, tím je jejich obraz roven množinám X,Y:

    \begin{align*}
      A_X &= \left(\begin{array}{ccc}
        3 & -2 & -1 \\
        2 & 0 & 0 \\
        0 & 1 & 0 \\
        0 & 0 & 2
      \end{array}\right) \\
      A_Y &= \left(\begin{array}{ccc}
        1 & 3 \\
        2 & -1 \\
        1 & 1 \\
        1 & 0
      \end{array}\right)
    \end{align*}

    Hledáme tedy množinu všech vektorů $\vec{p}$, které jsou vyjádřitelné jako lineární kombinace sloupců matice $A_X$, ale také $A_Y$.
    Nalezněme tedy všechny pětice $a,b,c,d,e$ splňující:

    \begin{align*}
      \vec{p} = \left(\begin{array}{ccc}
        3 & -2 & -1 \\
        2 & 0 & 0 \\
        0 & 1 & 0 \\
        0 & 0 & 2
      \end{array}\right) \cdot \left(\begin{array}{c}
        a\\b\\c
      \end{array}\right)
      &= \left(\begin{array}{ccc}
        1 & 3 \\
        2 & -1 \\
        1 & 1 \\
        1 & 0
      \end{array}\right) \cdot \left(\begin{array}{c}
        d\\e
      \end{array}\right) \\
      \left(\begin{array}{ccc}
        3 & -2 & -1 \\
        2 & 0 & 0 \\
        0 & 1 & 0 \\
        0 & 0 & 2
      \end{array}\right) \cdot \left(\begin{array}{c}
        a\\b\\c
      \end{array}\right) -
      \left(\begin{array}{ccc}
        1 & 3 \\
        2 & -1 \\
        1 & 1 \\
        1 & 0
      \end{array}\right) \cdot \left(\begin{array}{c}
        d\\e
      \end{array}\right) 
      &= 0 \\
      \left(\begin{array}{ccc}
        3 & -2 & -1 \\
        2 & 0 & 0 \\
        0 & 1 & 0 \\
        0 & 0 & 2
      \end{array}\right) \cdot \left(\begin{array}{c}
        a\\b\\c
      \end{array}\right) +
      \left(\begin{array}{ccc}
        -1 & -3 \\
        -2 & 1 \\
        -1 & -1 \\
        -1 & 0
      \end{array}\right) \cdot \left(\begin{array}{c}
        d\\e
      \end{array}\right) 
      &= 0 \\
      \left(\begin{array}{ccccc}
        3 & -2 & -1 & -1 & -3\\
        2 & 0 & 0 & -2 & 1\\
        0 & 1 & 0 & -1 & -1\\
        0 & 0 & 2 & -1 & 0
      \end{array}\right) \cdot \left(\begin{array}{c}
        a\\b\\c\\d\\e
      \end{array}\right)
      &= 0
    \end{align*}
  
  Vyřešme tuto soustavu lineárních rovnic:

  \[
    \left(\begin{array}{ccccc|c}
      3 & -2 & -1 & -1 & -3 & 0 \\
      2 & 0 & 0 & -2 & 1 & 0 \\
      0 & 1 & 0 & -1 & -1 & 0 \\
      0 & 0 & 2 & -1 & 0 & 0
    \end{array}\right)
    \overset{\text{Gauss--Jordan}}{\sim}
    \left(\begin{array}{ccccc|c}
      1 & 0 & 0 & 0 & \frac{27}{2} & 0 \\
      0 & 1 & 0 & 0 & 12 & 0 \\
      0 & 0 & 1 & 0 & \frac{13}{2} & 0 \\
      0 & 0 & 0 & 1 & 13 & 0
    \end{array}\right)
    \sim
    \left(\begin{array}{ccccc|c}
      2 & 0 & 0 & 0 & 27 & 0 \\
      0 & 1 & 0 & 0 & 12 & 0 \\
      0 & 0 & 2 & 0 & 13 & 0 \\
      0 & 0 & 0 & 1 & 13 & 0
    \end{array}\right)
  \]
  \[
    \overset{\text{Parametr } t\in\mathbb{R}}{\sim}
    \left(\begin{array}{ccccc|c}
      2 & 0 & 0 & 0 & 27 & 0 \\
      0 & 1 & 0 & 0 & 12 & 0 \\
      0 & 0 & 2 & 0 & 13 & 0 \\
      0 & 0 & 0 & 1 & 13 & 0 \\ \hline
      0 & 0 & 0 & 0 & 1 & -2t
    \end{array}\right)
    \sim
    \left(\begin{array}{ccccc|c}
      2 & 0 & 0 & 0 & 0 & 54t \\
      0 & 1 & 0 & 0 & 0 & 24t \\
      0 & 0 & 2 & 0 & 0 & 26t \\
      0 & 0 & 0 & 1 & 0 & 26t \\ \hline
      0 & 0 & 0 & 0 & 1 & -2t
    \end{array}\right)
    \sim
    \left(\begin{array}{ccccc|c}
      1 & 0 & 0 & 0 & 0 & 27t \\
      0 & 1 & 0 & 0 & 0 & 24t \\
      0 & 0 & 1 & 0 & 0 & 13t \\
      0 & 0 & 0 & 1 & 0 & 26t \\ \hline
      0 & 0 & 0 & 0 & 1 & -2t
    \end{array}\right)
  \]

  Dosad'me řešení do rovnice a zísejme všechny vektory $\vec{p}$:

  \begin{align*}
    \vec{p} &= \left(\begin{array}{ccc}
      3 & -2 & -1 \\
      2 & 0 & 0 \\
      0 & 1 & 0 \\
      0 & 0 & 2
    \end{array}\right) \cdot \left(\begin{array}{c}
      27t\\24t\\13t
    \end{array}\right)
    = \left(\begin{array}{ccc}
      1 & 3 \\
      2 & -1 \\
      1 & 1 \\
      1 & 0
    \end{array}\right) \cdot \left(\begin{array}{c}
      26t\\-2t
    \end{array}\right) \quad , t\in\mathbb{R} \\
    \vec{p} &= \left(\begin{array}{c}
      81t-48t-13t \\ 54t \\ 24t \\ 26t
    \end{array}\right)
    =  \left(\begin{array}{c}
      26t-6t \\ 52t + 2t \\ 26t-2t \\ 26t
    \end{array}\right) \quad , t\in\mathbb{R} \\
    \vec{p} &= \left(\begin{array}{c}
      20t \\ 54t \\ 24t \\ 26t
    \end{array}\right) \quad , t\in\mathbb{R}
  \end{align*}

  Pak tedy platí:

  \[
    X \cap Y = Span \left\{ \left(\begin{array}{c}
      20\\54\\24\\26
    \end{array}\right) \right\}
  \]

  Jednou z bází $X \cap Y$ je jednoprvková posloupnost vektorů $\left( (20,54,24,26)^T \right)$.

\end{enumerate}

\end{document}
